\section{Recursive routes, modes, and congestion}
\label{ch:routing}

\subsection{From paths to a recursion}

A path from \(i\) to \(j\) is a sequence of directed edges. If
\(\kappa_{kl}\) is an iceberg cost factor, a path cost is the product of its
edge costs. Equivalently, generalized costs add after taking logs. With a hard
minimum, the bilateral cost is
\[
 \widetilde\tau_{ij}
 =\min_{p\in\mathcal P^{ij}}\prod_{(k,l)\in p}\kappa_{kl}.
\]

The maintained model smooths this choice over feasible routes:
\[
 \tau_{ij}^{-\theta}
 =\mathbf 1\{i=j\}
  +\sum_{k\in\mathcal N^+(i)}
    \kappa_{ik}^{-\theta}\tau_{kj}^{-\theta}.
\]
For the economic-geography theorem, \(\theta=\sigma-1\). The urban model uses
its commuting heterogeneity parameter in the same transport equation. The
current package rejects an independent route curvature in the
economic-geography model because the paper's derivative is proved for the
common-curvature case.

The recursion can be written using a substochastic continuation matrix \(K\).
When its spectral radius is below one,
\[
 T=(I-K)^{-1}=I+K+K^2+\cdots.
\]
The resolvent sums all weighted walks. It also makes route derivatives
tractable: a change to one edge affects all origin-destination pairs whose
weighted walks use that edge.

\subsection{Recovering the baseline route object}

The package starts from directed edge traffic and location margins, not an
unobserved path table. It reconstructs the route kernel from the continuation
shares and obtains bilateral flows as
\[
 X^{OD}=\operatorname{diag}(o)\,
        (I-K)^{-1}\,
        \operatorname{diag}(a),
\]
where \(o\) is the source margin and \(a\) the destination absorption vector
in the relevant closure. For economic geography the source and destination
margins coincide with normalized income. For urban commuting they are
residence and workplace masses and can differ.

The reconstruction is accepted only if:
\begin{itemize}
  \item the route kernel is contractive;
  \item absorption probabilities add to one;
  \item bilateral rows and columns match the declared margins;
  \item reconstructed directed-edge traffic matches the input;
  \item every active policy edge has positive flow in its policy mode.
\end{itemize}

The fixed-route closure stores baseline origin-destination edge incidence.
The flexible-route closure differentiates the resolvent. Both are evaluated at
the same observed baseline, which makes their difference an exact local route
wedge rather than a comparison of two separately calibrated models.

\subsection{Choice-consistent modal aggregation}

Each directed edge can carry several modes. The active paper specification
uses
\[
 \kappa_e
 =\left(\sum_m\kappa_{e,m}^{-\eta}\right)^{-1/\eta},
 \qquad
 s_{e,m}
 =\frac{\kappa_{e,m}^{-\eta}}
        {\sum_n\kappa_{e,n}^{-\eta}}.
\]
The share is both the cost elasticity
\(\partial\log\kappa_e/\partial\log\kappa_{e,m}\) and the modal traffic share
under the maintained logsum. A lower modal cost raises its share. A larger
\(\eta\) makes modal allocation more responsive to relative costs.

The implementation represents the sign through a modal power:
\[
 p_m=
 \begin{cases}
 -\eta,&\text{\code{ChoiceLogsum}},\\
 +\eta,&\text{\code{ComponentCES}}.
 \end{cases}
\]
The positive-power alternative is retained for legacy replication only. It is
not interchangeable with the paper's choice model.

\subsection{Congestion maps}

Realized edge-mode cost can differ from the primitive infrastructure cost:
\[
 \log\kappa_{e,m}
 =\log\bar\kappa_{e,m}
  +\gamma_{e,m}\log\Xi_{e,m}
\]
for edge-local congestion. The headline application sets
\(\gamma_{e,\mathrm{road}}=0.092\) and other headline congestion elasticities to
zero.

The package also allows endpoint-terminal congestion. In that extension, an
edge-mode cost depends on total affected traffic at the declared origin and
destination terminals. This is useful for ports, rail terminals, or transit
hubs, but it is not part of the paper's main theorem application.

For a transport closure \(S\), collect congestion quantities in \(a\) and
write
\[
 x=Q_z^S z+C^S a,\qquad
 a=\theta+\Gamma^Sx.
\]
Eliminating \(a\) gives the transport response
\[
 (I-\Gamma^SC^S)^{-1}\Gamma^SQ_z^S.
\]
The code checks the transport matrix before adding its feedback to the spatial
Jacobian. It does not use ridge regularization when this system is singular.

\warningbox{Traffic must be measured in the units implied by the congestion
elasticity. A value-flow elasticity cannot be applied automatically to a
vehicle count, and a BPR elasticity cannot be inferred from link capacity
without a declared baseline volume-to-capacity ratio.}
